% 第 1 周周记（2026.08.03 - 08.07）
\documentclass[UTF8,fontset=fandol]{ctexart}
% =====================================================================
% internship/preamble.tex
% 实习周记统一模板的共用导言区。各周文件在 \documentclass 之后用
%   % =====================================================================
% internship/preamble.tex
% 实习周记统一模板的共用导言区。各周文件在 \documentclass 之后用
%   % =====================================================================
% internship/preamble.tex
% 实习周记统一模板的共用导言区。各周文件在 \documentclass 之后用
%   \input{../preamble.tex}
% 引入本文件，正文前用 \weekhead{周次}{时间}{主题} 输出抬头。
% 编译命令：xelatex weekN.tex
% =====================================================================

\usepackage[a4paper,top=2.6cm,bottom=2.6cm,left=2.8cm,right=2.8cm]{geometry}
\usepackage{booktabs}   % 表格横线，week3 的结果对比表会用到
\usepackage{array}

% 周记抬头：\weekhead{周次}{时间区间}{本周主题}
\newcommand{\weekhead}[3]{%
\begin{center}
{\zihao{3}\bfseries 实习周记（#1）}\\[10pt]
{\zihao{-4}姓名：【姓名】\qquad 学号：【学号】}\\[4pt]
{\zihao{-4}时间：#2}\\[4pt]
{\zihao{-4}主题：#3}
\end{center}
\vspace{2pt}
\hrule height 0.8pt
\vspace{10pt}
\zihao{-4}%
}

% 引入本文件，正文前用 \weekhead{周次}{时间}{主题} 输出抬头。
% 编译命令：xelatex weekN.tex
% =====================================================================

\usepackage[a4paper,top=2.6cm,bottom=2.6cm,left=2.8cm,right=2.8cm]{geometry}
\usepackage{booktabs}   % 表格横线，week3 的结果对比表会用到
\usepackage{array}

% 周记抬头：\weekhead{周次}{时间区间}{本周主题}
\newcommand{\weekhead}[3]{%
\begin{center}
{\zihao{3}\bfseries 实习周记（#1）}\\[10pt]
{\zihao{-4}姓名：【姓名】\qquad 学号：【学号】}\\[4pt]
{\zihao{-4}时间：#2}\\[4pt]
{\zihao{-4}主题：#3}
\end{center}
\vspace{2pt}
\hrule height 0.8pt
\vspace{10pt}
\zihao{-4}%
}

% 引入本文件，正文前用 \weekhead{周次}{时间}{主题} 输出抬头。
% 编译命令：xelatex weekN.tex
% =====================================================================

\usepackage[a4paper,top=2.6cm,bottom=2.6cm,left=2.8cm,right=2.8cm]{geometry}
\usepackage{booktabs}   % 表格横线，week3 的结果对比表会用到
\usepackage{array}

% 周记抬头：\weekhead{周次}{时间区间}{本周主题}
\newcommand{\weekhead}[3]{%
\begin{center}
{\zihao{3}\bfseries 实习周记（#1）}\\[10pt]
{\zihao{-4}姓名：【姓名】\qquad 学号：【学号】}\\[4pt]
{\zihao{-4}时间：#2}\\[4pt]
{\zihao{-4}主题：#3}
\end{center}
\vspace{2pt}
\hrule height 0.8pt
\vspace{10pt}
\zihao{-4}%
}


\begin{document}

\weekhead{第 1 周}{2026 年 8 月 3 日至 8 月 7 日}{课题入门：EVM 存储访问瓶颈与既有优化路线调研}

实习第一周，课题是"以太坊虚拟机存储访问混合预取方法研究"。我此前对以太坊的了解停留在共识和账户模型那一层，对执行层几乎没有什么概念，所以本周的主要任务是补背景，把"存储访问到底慢在哪里"这个问题弄清楚。导师给了一份阅读路线，先读 Erigon 的块内状态缓存实现，再读预取和 learned index 方向的综述，我自己又补了以太坊文档和几个 EIP。材料读到一半时有种感觉，性能问题的讨论如果离了延迟数字，很容易变成空对空。

印象最深的是几个具体数字。EVM 指令里，storage 读写是唯一需要数据库 I/O 的操作，其它指令基本都在寄存器、内存这些纳秒级的地方完成。EIP-7707 提到，数据加载在某些客户端里能占到出块执行时间的七成以上。Erigon 用块内缓存把同一块内的重复访问变成热访问，可每个槽位在块内的第一次访问仍然要读数据库，这部分绕不过去。导师让我在本机把这个差距复现出来，我用 Go 写了一个小基准：热访问用 map 模拟，冷访问用 O_DIRECT 绕过页缓存直接读文件来模拟，结果热访问约 30 纳秒，冷访问约 3 微秒，差了约 100 倍。数值本身并不新奇，但自己亲手跑出来之后，对课题动机的理解确实不一样了。

文献调研沿着三个方向展开：EVM 执行优化、数据预取、learned index。EVM 方向从存储引擎、access list、静态分析到并行执行都有不少工作，但没有人做预测式的预取；数据预取方向的主流做法是把 ML 放在关键路径上在线推理；learned index 里 Shift-Table 那种"模型给粗位置、算法层再修正"的混合结构最让我受启发。我也翻了 Erigon 的源码，几十万行规模的 Go 工程，想通过直接修改客户端来验证想法，迭代周期会非常长，这个判断影响了后面整个实验路线的选择。

周五和导师过了一遍调研结论，讨论最后落到一个具体问题上：预取到底有没有收益空间，收益有多大。光靠推理和读文献回答不了，需要真实的主网数据。所以下周的计划是先把数据采集环境搭起来，把交易 trace 拿到手。第一周最大的感受是，一个性能课题的起点，常常是一组说得清来源的延迟数字，外加一个能重复的测量方法。

\end{document}
